\chapter{Decomposition}
\label{ch:decomposition}

\section*{What this chapter is for}

How to cut a build into pieces small enough that each can be handed to somebody
\dash{} or something \dash{} that never saw the conversation. It is the skill
that makes parallel work possible, it is the skill that assisted building
actually requires, and it is measurable in advance.

\section{The atomicity test}

A task is atomic when all five hold. Fewer than five and it will come back
wrong or block something else.

\begin{kitmethod}
\textbf{1. One area of the problem.} Not one layer \dash{} one area. A task
spanning two areas has two reasons to change and two people with opinions.

\textbf{2. It owns its files.} No file is written by two tasks running at the
same time. This is the condition that creates parallelism and the one most
often violated.

\textbf{3. One acceptance criterion, with a check.} If you cannot say what
proves it done, it is not a task.

\textbf{4. Green in, green out.} It starts from a working state and ends in
one. A task that leaves the build broken cannot be handed to anybody.

\textbf{5. Self-sufficient from its own text.} Everything needed is in the
task description. ``From the conversation earlier'' is a defect, not an
inconvenience \dash{} it means the task cannot be given to a fresh session, and
a fresh session is exactly what you will give it to.
\end{kitmethod}

\section{File ownership is the whole mechanism}

Condition two deserves its own section because everything about parallel work
follows from it.

Three files in almost any project want to be touched by every task: the
composition root or entry point, the data-context or schema definition, and the
migration or configuration directory. If two tasks both add a line to the entry
point, they conflict, and the conflict is guaranteed rather than unlucky.

\begin{kitmethod}
\textbf{Those files belong to the join, not to any task.}

Practically: declare everything they need \emph{up front}, in one task, before
any parallel work starts. Every entity, every registration, every route group.
It costs a few extra minutes in the first task and removes every conflict
afterwards.

A parallel task never edits a shared file. It exposes something in its own file
\dash{} a function, a module, a registration helper \dash{} and you add the one
line at the join.
\end{kitmethod}

\judgement{This is the single most useful thing in this chapter for assisted
building, because a model given two tasks that both touch the entry point will
cheerfully produce two incompatible versions of it and neither will look
wrong.}

\section{Width and depth}

Two numbers, computed before you start, that tell you whether the build fits.

\textbf{Width} is the sum of all the task estimates: total work.

\textbf{Depth} is the longest chain of tasks that must happen in order: the
critical path.

Depth is the one that binds. A build that is 140 minutes wide and 35 minutes
deep fits in an hour if you can run things in parallel, and does not fit at all
if you cannot. A build that is 60 minutes wide and 55 deep does not benefit
from any parallelism at all.

\begin{kitnote}
Computing these takes five minutes with the task list in front of you, and it
converts ``I think this fits'' into a statement with two numbers in it. It also
tells you which task to start first: the one at the head of the longest chain,
regardless of how appealing the others look.
\end{kitnote}

\section{What you never hand off}

Some work is yours whatever the assistance available.

\begin{itemize}
  \item Discovery. It is the thing being assessed.
  \item The design. Same reason.
  \item The hardest slice \dash{} the conflict case, the irreversible
        operation. Chapter~\ref{ch:ai-visibly} covers why.
  \item Its check.
  \item The demonstration.
\end{itemize}

\judgement{Notice that this list is close to the published assessment criteria
for a build round. That is not a coincidence: the work you must not delegate is
the work being assessed, and the work that is safe to delegate is the work that
is safe precisely because nobody is grading it.}

\section{The task text}

A task handed to a fresh session needs, in its own words: what to build, which
files it owns, what proves it done, and what it may assume already exists.

Three things it must \emph{not} contain: a copy of your general working
procedure, which will drift from the original; absolute paths from your
machine; and any credential, ever \dash{} name the variable, never its value.

\begin{drill}
Take a small system you could build in two hours and decompose it on paper:
tasks, owned files, dependencies. Compute width and depth. Then check condition
two across the whole list and count the violations \dash{} there are almost
always two or three, and they are almost always the entry point and the schema.
\end{drill}

\section*{What this chapter does not claim}

None of this is measured. It is a working method, and the width-and-depth
calculation in particular is a rough instrument \dash{} estimates at this
granularity are not accurate, and the value is in the ratio rather than in
either number.

It does not claim parallel work is always worth it. Below a certain size the
coordination costs more than it saves, and in a session short enough the
correct number of parallel tasks is zero.
